\section*{Capítulo \ref{ch:g4:clouds-rain-snow} --- Nuvens, chuva e neve}

\begin{solution}{exo:g4:clouds-rain-snow:1}
Uma \omterm{def:g4:clouds-rain-snow:cloud}{nuvem} é feita de uma enorme quantidade de gotículas de água, ou de
pedacinhos minúsculos de gelo, flutuando no ar. A \omterm{def:g4:clouds-rain-snow:cloud}{neblina} é uma \omterm{def:g4:clouds-rain-snow:cloud}{nuvem} que
encosta no chão.
\end{solution}

\begin{solution}{exo:g4:clouds-rain-snow:2}
(a) \omterm{def:g4:clouds-rain-snow:cloud-families}{Estrato}. (b) \omterm{def:g4:clouds-rain-snow:cloud-families}{Cirro}. (c) \omterm{def:g4:clouds-rain-snow:cloud-families}{Cúmulo}. (d) \omterm{def:g4:clouds-rain-snow:cloud-families}{Cumulonimbo}.
\end{solution}

\begin{solution}{exo:g4:clouds-rain-snow:3}
(a) \omterm{def:g4:clouds-rain-snow:precipitation}{Granizo}. (b) \omterm{def:g4:clouds-rain-snow:precipitation}{Chuva}. (c) \omterm{def:g4:clouds-rain-snow:precipitation}{Neve}.
\end{solution}

\begin{solution}{exo:g4:clouds-rain-snow:4}
(a) Um \omterm{def:g4:clouds-rain-snow:rain-gauge}{pluviômetro}. (b) Um termômetro. (c) Um \omterm{def:g4:clouds-rain-snow:wind-vane}{cata-vento}.
\end{solution}

\begin{solution}{exo:g4:clouds-rain-snow:5}
Caíram 17 milímetros de \omterm{def:g4:clouds-rain-snow:precipitation}{chuva}. $20 - 17 = 3$: mais 3 milímetros levariam
a água até a marca 20.
\end{solution}

\begin{solution}{exo:g4:clouds-rain-snow:6}
O dia mais chuvoso foi a quarta-feira (12 milímetros). Não choveu na
segunda nem na sexta. A quinta-feira teve 7 milímetros.
\end{solution}

\begin{solution}{exo:g4:clouds-rain-snow:7}
Os flocos de \omterm{def:g4:clouds-rain-snow:precipitation}{neve} se formam no alto da \omterm{def:g4:clouds-rain-snow:cloud}{nuvem} fria. Lá em cima, na vila da
\omterm{def:g1:landscapes-around-us:landform}{montanha}, o ar está abaixo de zero até embaixo; por isso eles chegam como
\omterm{def:g4:clouds-rain-snow:precipitation}{neve}. Mais embaixo, no \omterm{def:g1:landscapes-around-us:landform}{vale}, o ar está acima de \SI{0}{\celsius}: ao
atravessá-lo na descida, os flocos derretem e chegam ao chão como \omterm{def:g4:clouds-rain-snow:precipitation}{chuva}.
\end{solution}

\begin{solution}{exo:g4:clouds-rain-snow:8}
O gráfico de barras tem sete barras de alturas 3, 0, 1, 9, 6, 0 e 1 (sem
barra na terça e no sábado). Na semana caíram
$3 + 0 + 1 + 9 + 6 + 0 + 1 = 20$ milímetros. A quinta-feira trouxe 9 dos
20 milímetros: $\frac{9}{20}$ da \omterm{def:g4:clouds-rain-snow:precipitation}{chuva} da semana.
\end{solution}

\begin{solution}{exo:g4:clouds-rain-snow:9}
Os muros e as árvores protegem da \omterm{def:g4:clouds-rain-snow:precipitation}{chuva} o chão perto deles, principalmente
quando o vento joga a \omterm{def:g4:clouds-rain-snow:precipitation}{chuva} de lado, e gotas podem pingar das folhas
dentro do \omterm{def:g4:clouds-rain-snow:rain-gauge}{pluviômetro}: perto deles, o \omterm{def:g4:clouds-rain-snow:rain-gauge}{pluviômetro} recolheria \omterm{def:g4:clouds-rain-snow:precipitation}{chuva} de
menos, ou de mais. O termômetro tem de medir a temperatura do ar; em pleno
sol, ele seria aquecido pelo próprio sol e marcaria uma temperatura alta
demais. O abrigo branco de venezianas impede a entrada do sol e deixa o ar passar.
\end{solution}

\begin{solution}{exo:g4:clouds-rain-snow:10}
Um \omterm{def:g4:clouds-rain-snow:wind-vane}{cata-vento} aponta para onde o vento vem: o vento vem do oeste. Ele está
indo para o leste.
\end{solution}

\begin{solution}{exo:g4:clouds-rain-snow:11}
Um milímetro de \omterm{def:g4:clouds-rain-snow:precipitation}{chuva} sobre um metro quadrado é um litro. Sobre um metro
quadrado, 8 milímetros dão 8 litros. Sobre 500 metros quadrados:
$500 \times 8 = 4000$ litros. Com regadores de 10 litros:
$4000 \div 10 = 400$ regadores.
\end{solution}
